\documentclass[xcolor=svgnames]{beamer}
\usepackage{latexsym,amssymb,amsmath,stmaryrd}
\usepackage[utf8]{inputenc}
\usepackage[T1]{fontenc}
\usepackage[francais]{babel}
\usepackage{pgf}
\usepackage{tikz}
\usetikzlibrary{arrows,patterns,plotmarks,shapes,snakes,er,3d,automata,%
backgrounds,topaths,trees,petri,mindmap}
\usepackage{pgfbaseimage}
\usepackage{ifpdf}
\ifpdf
\usepackage{graphicx}
\else
%\usepackage[dvips]{graphicx}
\usepackage{pstricks,pst-tree,pst-node}
\newpsobject{showgrid}{psgrid}{subgriddiv=1,griddots=10,gridlabels=6pt}
\fi
\usepackage{moreverb}
\usepackage[normalem]{ulem}
\usepackage{version}
\usepackage{url}
\usepackage{multicol}
\usepackage{listings}
\lstset{
  language=[ANSI]C, 
%  gobble=2, 
%  escapeinside="", 
  basicstyle=\ttfamily, 
%  directivestyle=\color{Fuchsia},
  identifierstyle = \color{DarkOrange},
  keywordstyle=\color{DarkGreen}, 
  commentstyle=\color{red}, 
%  numbers=left, 
%  numbersep=5pt, 
%  numberstyle=\scriptsize, 
%  moredelim=[is][\color{gray}\itshape]{/*}{*/}, 
%  moredelim=[is][\alert]{/+}{+/}, 
%  morecomment=[is]{/=}{=/}, 
%  morecomment=[in]{comment=}{,} 
%  fancyvrb=true 
} 

\newcommand{\binaire}[1]{\ensuremath{\underline{#1}}}
\newcommand{\C}[1]{\texttt{#1}}
\newcommand{\bbbn}{\ensuremath{\mathbb{N}}}
\input{communs/cc_beamer}
\usetheme{classic}
\newcommand{\nowrite}{\put(10,-4){\includegraphics[scale=.05]{creative_commons/nopencil}}}
\newcommand{\youwrite}{\put(10,-4){\includegraphics[scale=.05]{creative_commons/pencil}}}
\newcommand{\aemporter}{\put(10,-6){\includegraphics[scale=.05]{creative_commons/szymonraj_Shopping_bag}}}


%%% Titre -- Cours 1
\title{Bases de programmation -- Cours 1. Introduction\\ L'instruction de contrôle \C{if}}

% Éléments d'architecture des
% ordinateurs, mini-assembleur\\La programmation
%   structurée en langage C\\ L'instruction de contrôle \C{if} premier programme C, if else}
\author{Pierre Boudes}
\date{\today}

\begin{document}

%% Page de titre et licence CC.
\begin{frame}
	\titlepage
	\vfill
	\begin{center}
		\CcGroupByNcSa{0.83}{0.95ex}\\[2.5ex]
		{\tiny\CcNote{\CcLongnameByNcSa}}
		\vspace*{-2.5ex}
	\end{center}
\end{frame}

\section[Langages et programmes]{Introduction : langages et programmes}
\begin{frame}
  \frametitle{Langages et programmes}

\begin{columns}
\column[T]{4.5cm}
\structure{Humains}\\
\only<2-2>{
\alert{communiquent entre eux en langage naturel}
}
\uncover<3->{
\alert{langage naturel}
}
\column[T]{3.5cm}
\uncover<7->{\alert{Langage de programmation}}
\begin{tikzpicture}
\useasboundingbox (0,0) rectangle (2,4.3);
\begin{scope}[shift={(6.5,3)}]%computer 3
\draw[ultra thick, rounded corners=4pt,fill=blue!15] (0,0) rectangle (1,0.7);
\draw[ultra thick, pattern=crosshatch dots] (0,0) -- (-0.4,-0.4) -- (0.6,
-0.4) -- (1, 0) -- cycle;
\end{scope}
\begin{scope}[shift={(5,2.5)}]%computer 1
\draw[ultra thick, rounded corners=4pt,fill=blue!10] (0,0) rectangle (1,0.7);
\draw[ultra thick, pattern=crosshatch dots] (0,0) -- (-0.4,-0.4) -- (0.6,
-0.4) -- (1, 0) -- cycle;
\end{scope}
\begin{scope}[shift={(5.6,1.9)}]%computer 2
\draw[ultra thick, rounded corners=4pt,fill=blue!20] (0,0) rectangle (1,0.7);
\draw[ultra thick, pattern=crosshatch dots] (0,0) -- (-0.4,-0.4) -- (0.6,
-0.4) -- (1, 0) -- cycle;
\end{scope}
\begin{scope}[shift={(-3,2.5)}]%hum 1
  \draw[ultra thick,fill=pink] (0,0.9) circle (0.3cm); \draw[draw,line
  width=1.1mm,fill=pink, rotate=45, rounded corners=4pt] (-0.7, -0.15)
  rectangle (0.7, 0.15) (-0.15, -0.7) rectangle (0.15, 0.7) ;
  \draw[draw=none,fill=blue, rotate=45, rounded corners=4pt] (-0.7,
  -0.15) rectangle (0.7, 0.15) (-0.15, -0.7) rectangle (0.15, 0.7) ;
\end{scope}
\begin{scope}[shift={(-4,1)}]%hum 2
  \draw[ultra thick,fill=LightSalmon] (0,0.9) circle (0.3cm); \draw[draw,line
  width=1.1mm,fill=pink, rotate=45, rounded corners=4pt] (-0.7, -0.15)
  rectangle (0.7, 0.15) (-0.15, -0.7) rectangle (0.15, 0.7) ;
  \draw[draw=none,fill=DarkTurquoise, rotate=45, rounded corners=4pt] (-0.7,
  -0.15) rectangle (0.7, 0.15) (-0.15, -0.7) rectangle (0.15, 0.7) ;
\end{scope}
\begin{scope}[shift={(-2,1)}]%hum 3
  \draw[ultra thick,fill=DarkRed] (0,0.9) circle (0.3cm); \draw[draw,line
  width=1.1mm,fill=pink, rotate=45, rounded corners=4pt] (-0.7, -0.15)
  rectangle (0.7, 0.15) (-0.15, -0.7) rectangle (0.15, 0.7) ;
  \draw[draw=none,fill=DarkGreen, rotate=45, rounded corners=4pt] (-0.7,
  -0.15) rectangle (0.7, 0.15) (-0.15, -0.7) rectangle (0.15, 0.7) ;
\end{scope}
\only<3-6>{
\draw[draw,line width=4mm,->,yellow](-2,3)--node [black] (4,3) {utilisent} (4,3);
}
\only<4-6>{%programment
\draw[draw,line width=4mm,->,Goldenrod](-1.9,2.5)--node [black] (4,3) {programment} (4.1,2.5);
}
\only<5-6>{%binaire
\draw [draw=none,text=blue] (4.7,3.5) node {01101} -- +(-0.2,-0.4)
node {10100};
\draw [draw=none,text=blue] (7,1.5) node {10101} -- +(-0.5,-0.4) node
{00101};
\draw [draw=none,text=red] (-2,3.5) node {!!?} -- +(-0.2,-0.4)
node {??};
}
\only<6-6>{%asm
\draw [draw=none,text=DarkGreen] (5.1,4.1) node {\C{movb 0x65, \%eax}};
\draw [draw=none,text=DarkGreen] (4.1,1.2) node {\C{movl \%eax, \%edx}};
}

\only<7->{%source
\draw[draw,line width=4mm,->,Goldenrod](-1.9,2.5)--node [black]
(0.4,3) {programment} (0.3,2.5);
\node [rectangle, draw, fill=blue!20,
text width=3em, text centered, minimum height=4.6em] at (1.2,2.5) {};
\node [rectangle, draw, fill=blue!20,
    text width=3em, text centered, minimum height=4.6em] at (1.1,2.6)
    {};
\node[line width=0mm,Goldenrod,rotate=60] at (1.1,2.6) {\Large\textbf{source}};
\node at (1.1,1.9) {\C{prog.c}};
}
\only<8->{%traduction prise en charge par un programme
\begin{scope}[shift={(2.4,2.2)},black!30]%computer 4
\draw[ultra thick, rounded corners=4pt,fill=blue!10] (0,0) rectangle (1,0.7);
\draw[ultra thick, pattern=crosshatch dots] (0,0) -- (-0.4,-0.4) -- (0.6,
-0.4) -- (1, 0) -- cycle;
\end{scope}
\draw[draw,line width=4mm,->,blue!20](2,2.5)--node [black] (4,3)
{traduction} (4.1,2.5);
}
\only<9->{%algorithme
\node [draw, text badly ragged, cloud,cloud ignores aspect, cloud
puffs=11, fill=green!20,text width=2cm]
at (-3.1,3) {algorithmes\\ \alert{(pseudocode)}};
}
\end{tikzpicture}
\column[T]{4.5cm}
\structure{Ordinateurs}\\
\only<2-2>{
\alert{obéissent à des instructions en langage machine}
}
\uncover<3->{
\alert{langage machine}
}
\end{columns}
\uncover<10->{
Dans ce cours : surtour du \textbf{langage C}, du pseudocode pour mettre au
point les algorithmes et un petit peu
d'un langage assembleur \emph{jouet} (amil) pour expliquer le
langage C.
}
\end{frame}

\mode<beamer>{\section[Plan]{}
\frame[label=plan]{\tableofcontents}}


\begin{frame}
  \begin{quote}
  « L'informatique n'est pas plus la science des ordinateurs que
  l'astronomie n'est celle des télescopes. »\\
  \hfill E.~W.~Dijkstra\\[1cm]
\end{quote}
\end{frame}

\section[Langage machine]{Architecture matérielle et langage machine}
\subsection{Architecture de von Neumann}
\begin{frame}
  \frametitle{Architecture de von Neumann\nowrite}
  \begin{itemize}
  \item John William Mauchly et John Eckert autant (ou plus) que
    vN\pause
  \item Qu'est-ce que c'est ?\pause
    \begin{itemize}
    \item L'idée d'une machine à \alert{programme
        stocké}\pause
    \item  Une machine réalisée, l'ancêtre de tous nos processeurs\pause
    \end{itemize}
  \item De quoi  cette machine est-elle faite ?\pause
  \begin{columns}[t]
  \begin{column}{4cm}
\centering
\structure{Machine de vN}

  \begin{tikzpicture}[draw=gray, line width=2pt]
    \useasboundingbox (.9,0) rectangle (4,2.5);
    \draw[very thin] (-.2,-1.9) rectangle (5.5,2.48);
    % memoire
    \node at (4.5,2.3) {mémoire};
    \draw[top color=blue!20] (4,-1.8) rectangle (5,2);
    \draw[loosely dashed, white](3.95,-1.8) -- (5.05,-1.8);
%    \draw[draw=none,fill=blue!20] (4,-2) rectangle (5,2);
    \foreach \y in {1,2,...,7} {
      \draw[draw=blue!40, line width=1pt] (4,2 - \y/2) -- (5, 2 -
      \y/2);
      \node[red] at (5.2, 2 - \y/2 + 0.25) {\y};
    }
\only<7->{%uc
    \node at (1,1.65) {UC};
    \draw[rounded corners=4pt] (0, -1.35) rectangle (2,1.35);
    \foreach \y in {-1.3,-0.7,...,1} {
      \draw[left color=red!20, right color=blue!20] (0.125,\y + 0.2)
      rectangle (0.925,\y + 0.6);
      \draw[left color=red!20, right color=blue!20] (1.075,\y + 0.2)
      rectangle (1.875,\y + 0.6);
    }
}
    \only<8->{%bus
    \draw[draw=none,fill=gray] (2,-0.6) rectangle (4,0.6);
    \draw[draw=none,fill=blue!20] (2,-0.5) rectangle (4,-0.05);
%    \draw[<->,blue!50] (2,-0.25) -- (4,-0.25); trop chargé.
    \draw[draw=none,fill=red!20] (2,0.05) rectangle (4,0.5);
    \node at (3,0.9) {bus};
    \node[red!40!black!60] at (3, 0.3) {adresses};
    \node[blue!40!black!60] at (3, -0.25) {données};
}
  \end{tikzpicture}
  \end{column}
  \begin{column}{6.5cm}
\vspace{-0.5cm}
    \begin{itemize}
    \item De mémoire (une suite de cases numérotées)\pause
    \item d'une \alert{unité de calcul}, travaillant sur des \alert{registres}\pause
    \item d'un \alert{bus} système (adresses et données) reliant mémoire
      et UC\pause
    \item \sout{De périphériques} (on oublie !)\pause
    \item La mémoire contient le programme et les données.
    \end{itemize}
\end{column}
\end{columns}
  \end{itemize}
\end{frame}
%\end{document}

\subsection{Représentation  des informations}

\begin{frame}
  \frametitle{Représentation en binaire des informations}
%  \pause
  \begin{definition}[bit]
    \begin{itemize}
    \item Le chiffre binaire, ou \emph{bit}, est l'équivalent binaire
      de nos chiffres décimaux. Il peut valoir soit $0$ soit $1$. Un
      bit est une \alert{quantité élémentaire d'information} (oui ou non,
      ouvert ou fermé, etc.).\pause
    \item L'information manipulée par un ordinateur est faite de
      bits.\pause
\item Les cases mémoires et les registres contiennent des \alert{mots mémoire}
  : des suite de $n$ bits, où $n$ est fixé une fois pour toute par
  l'architecture matérielle.\pause
\item les instructions du langage machine sont écrites en binaire.\pause
\item   le \alert{langage assembleur} est une notation du langage
  machine plus pratique pour les humains.
   \end{itemize}
  \end{definition}
\pause Vous en verrez plus sur les codages en binaire des données
dans un autre cours.
\end{frame}

\subsection{Exécution des instructions}
\begin{frame}
  \frametitle{Cycle d'exécution\nowrite}
\pause
\begin{itemize}
  \item Le registre \alert{compteur de programme} (CP) contient l'adresse du mot mémoire
    représentant la prochaine instruction\pause
\item le contenu de ce mot est transféré de la mémoire centrale dans
  le \alert{registre d'instruction} (RI)\pause
\item CP est \emph{incrémenté} (c'est à dire que sa valeur augmente de 1)\pause
\item le contenu de RI est décodé afin de déterminer l'opération à
  exécuter\pause
\item l'opération est exécutée (le contenu d'un ou plusieurs
  registres est modifié, ou bien celui d'une case mémoire)\pause
\item Fin du cycle d'exécution et démarrage d'un nouveau cycle
\end{itemize}
\end{frame}

\begin{frame}
  \frametitle{Instructions\nowrite}
Une instruction machine  type comporte un \alert{code d'opération} et, si nécessaire, une
ou deux \emph{opérandes} (ou \emph{arguments} de l'opération).

\pause
\begin{block}{Vocabulaire} Dans l'expression arithmétique usuelle $3 + 5$,  le signe $+$ est l'opérateur
et les nombres $3$ et $5$ sont les opérandes.
\end{block}
\end{frame}

\begin{frame}
\frametitle{Quelques instructions typiques (Amil)\nowrite}
\pause
  \begin{tabular}[c]{lp{8.3cm}}
  \C{stop} & Arrête l'exécution du programme.\\
  \C{noop} & N'effectue aucune opération.\\ \pause
  \C{lecture i rj} & Charge, dans le registre $j$, le contenu de la mémoire d'adresse $i$.\\
  \C{ecriture ri j} & Écrit le contenu du registre $i$ dans la mémoire
  d'adresse $j$.\\ \pause
  \C{saut i} & Met CP à la valeur $i$.\\
  \C{sautpos ri j} & Si la valeur contenue dans le registre $i$ est
  positive ou nulle, met CP à la valeur $j$.\\ \pause
%  \C{valeur x ri} & Initialise le registre $i$ avec la valeur $x$.\\
 \C{inverse ri} & Inverse le signe du contenu du registre $i$.\\
  \C{add ri rj} & Ajoute la valeur du registre $i$ à celle du registre $j$.\\
  \C{soustr ri rj} & Soustrait la valeur du registre $i$ à celle du registre $j$.\\
  \C{mult ri rj} & Multiplie \ldots\\
  \C{div ri rj} & Divise \ldots \\ \pause
% \\
%   \multicolumn{2}{l}{\textbf{Instructions plus avancées}}\\
%   \C{et ri rj} & Effectue le et bit à bit de la valeur du registre
%   $i$ et de celle du registre $j$. Le résultat est placé dans le
%   registre $j$.\\
   \C{lecture *ri rj} & Charge, dans $rj$, le contenu de la mémoire dont l'adresse est la valeur du registre $i$\\
%   \C{ecriture ri *rj} & Écrit le contenu de $ri$ dans la  mémoire dont  l'adresse est la valeur du registre $j$.
\end{tabular}
\end{frame}


%\subsection{Trace d'exécution}
\begin{frame}[fragile]
\frametitle{Trace d'exécution\nowrite}
On simule pas à pas l'exécution du programme.
\begin{columns}[t]
  \begin{column}{3.5cm}
\newcommand{\N}[1]{\structure{\scriptsize #1.}}
\pause \structure{Programme}
\begin{semiverbatim}
\N{ 1} lecture 10 r0
\N{ 2} lecture 11 r2
\N{ 3} soustr r2 r0
\N{ 4} sautpos r0 8
\N{ 5} lecture 10 r2
\N{ 6} add r2 r0
\N{ 7} saut 4
\N{ 8} ecriture r0 12
\N{ 9} stop
\N{ 10} 14
\N{ 11} 5
\N{ 12} ?
\end{semiverbatim}
 \end{column}
  \begin{column}{10cm}
    \pause \structure{Trace}\\[0.5cm]
\scriptsize
\newcommand{\commentaire}[1]{\pause} %hack
\begin{tabular}[c]{l||c|c|c|c|c|c|c|}
\hline
 \emph{Instructions} & Cycles & CP& r0& r2& 10& 11 &12\\ \hline \pause
\hfill Initialisation & 0 & 1 & ? & ? & 14 & 5 & ?
 \\ \hline \commentaire{Lecture de la donnée d'adresse 10 dans le
   registre 0}
\C{lecture 10 r0} & 1 & 2  & 14 & & & &\\ \hline
 \commentaire{Initialisation du registre 2 à 5
} \C{lecture 11 r2} & 2 & 3  & & 5 & & &\\ \hline
 \commentaire{Soustrait la valeur du registre 2 au registre 0
} \C{soustr r2 r0
} & 3 & 4  & 9 & & & &\\ \hline
 \commentaire{Si la valeur (9) du registre 0 est positive, saute à l'adresse 8
} \C{sautpos r0 8
} & 4 & \textbf{8} & & & & &\\ \hline
 \commentaire{Écriture du registre 0 à l'adresse 12
} \C{ecriture r0 12
} & 5 & 9  & & & & & 9
\\ \hline
 \commentaire{Fin du processus.
} \C{stop
} & 6 & 10  & & & & &\\ \hline
\end{tabular}
  \end{column}
\end{columns}
\end{frame}


\section[Le langage C]{Premiers pas en langage C}
\subsection{La programmation structurée}
\begin{frame}
  \frametitle{La programmation structurée}

  \begin{definition}[Programmation structurée]
    Programmer par \structure{blocs} d'instructions en combinant ces blocs de trois manières :\pause
    \begin{enumerate}
    \item exécuter les blocs les uns à la suite des autres (\structure{séquence})\pause
    \item si une certaine condition est vraie, exécuter un bloc sinon
      en exécuter un autre (\structure{sélection})\pause
    \item recommencer l'exécution d'un bloc tant qu'une certaine
      condition est vraie (\structure{répétition}).\pause
    \end{enumerate}
    Un bloc peut lui-même être une combinaison de blocs (ou juste une instruction).
  \end{definition}
\pause La sélection et la répétition sont assurées par des
\structure{instructions de contrôle}.
  \pause Tout programme en langage machine peut
  être transformé en un programme structuré (Böhm-Jacopini 1966).
\end{frame}


\begin{frame}
  \frametitle{Nécessité d'une traduction\nowrite}
  Les langages structurés, dits de haut niveau, nécessitent une \alert{traduction
  en langage machine}.\pause

Cette traduction est assurée par un programme particulier : un compilateur
 (traduction une fois pour toute) ou un interprète (traduction à
 chaque exécution).\pause

 La traduction suit des règles précises et systématiques. À chaque
 instruction du langage correspond un \structure{schéma de
   traduction}. C'est ce schéma qui donne sont sens (son effet) à
 l'instruction.
\end{frame}
%\subsection{Structure d'un programme C}
\begin{frame}[fragile]
  \frametitle{Construction d'un programme C\nowrite}
\structure{Code source}
\begin{lstlisting}[basicstyle=\ttfamily\scriptsize]
/* Declaration de fonctionnalites supplementaires */
#include <stdlib.h> /* EXIT_SUCCESS */

/* Declaration des constantes et types utilisateurs */

/* Declaration des fonctions utilisateurs */

/* Fonction principale */
int main()
{
    /* Declaration et initialisation des variables */
    ...
    /* valeur fonction */
    return EXIT_SUCCESS;
}

/* Definitions des fonctions utilisateurs */
\end{lstlisting}

Les commentaires sont ignorés lors de la traduction en langage machine.
\end{frame}

\subsection{Variables impératives et affectation}
\begin{frame}[fragile]
  \frametitle{Traduction de l'affectation}
  \begin{columns}
 \begin{column}[t]{2.5cm}
\structure{Code source}
\begin{lstlisting}[basicstyle=\ttfamily\scriptsize]
...
/* Fonction principale */
int main()
{
    /* Declaration et initialisation des variables */
    int x = 5;



    x = 3 * x + 1;
    ...
\end{lstlisting}
\end{column}
\pause
\begin{column}[t]{6.5cm}
\hfill\structure{Schéma de traduction}\\[.3cm]
\tikzstyle{amil}=[rectangle, draw,  fill=blue!20, text badly ragged,
text width=5.4cm]
\tikzstyle{line} = [draw, thick, -latex']
\begin{tikzpicture}[node distance=1.5cm]
\useasboundingbox (-4,-2) rectangle (3,1.5);
\uncover<2-2> {%
\node [amil,cloud, cloud ignores aspect, cloud puffs = 15, fill=green!20] (avant) {Les déclarations ne sont pas traduites en
    instructions, mais en réservations mémoires};
}
\uncover<3-> {%
\node [amil,fill=green!20] (avant) {\emph{$x$ est une case mémoire de
    type entier qui contient initialement 5}};
\draw [draw, thick, green] (-7.1,-0.5) rectangle (-5.2, -0.1) --
(-5.2,-0.3); \draw [ thick, green, ->] (-5.2,-0.3)  -> (avant);
}
\uncover<4-4>{%
  \node [amil,cloud, cloud ignores aspect, cloud puffs = 15, below of=avant,node distance=2cm] (eval) {l'affectation est traduite par
    l'évaluation d'une expression et  une écriture en mémoire};
}
\uncover<5->{%
  \node [amil,below of=avant] (eval) {\emph{évaluation de l'expression
    $3 x + 1$ dans un registre}};
\draw [draw, thick, blue] (-6.4,-1.84) rectangle (-4.7, -1.44) --
(-4.7,-1.64) ; \draw[ thick, blue, ->] (-4.7,-1.64)  -- (eval.west);
}
\uncover<6->{%
  \node [amil,below of=eval,node distance=1.1cm] (eval) {\emph{écriture du résultat à l'adresse
  de $x$}};
\draw [draw, thick, blue] (-7.1,-1.84) rectangle (-6.5, -1.44);
\draw  [thick, blue, ->]   (-6.8,-1.84) -- (eval.west);
}
\end{tikzpicture}
\end{column}
\end{columns}
\end{frame}


\subsection{L'instruction de contrôle if}
\begin{frame}[fragile]
  \frametitle{L'instruction de contrôle if}

\structure{Syntaxe :} \C{if (\emph{condition}) \{\emph{\colorbox{green!20}{bloc1}}\} else
\{\emph{\colorbox{red!20}{bloc2}}\}}.
\medskip

\pause
  \begin{columns}
\begin{column}[t]{5cm}
\structure{Code source}
\begin{verbatim}
/* avant */
if (age < 18)
{
  permis = 0;
}
else
{
  permis = 1;
}
/* après */
\end{verbatim}
   \end{column}
\pause
   \begin{column}[t]{4cm}
\structure{Schéma de traduction}\\[.3cm]
\tikzstyle{amil}=[rectangle, draw,  fill=blue!20, text badly ragged, text width=3.7cm]

\tikzstyle{line} = [draw, thick, -latex']
\begin{tikzpicture}[node distance=.7cm]
  \node [amil] (avant) {\emph{avant}};
  \node [amil,below of=avant] (eval) {évalue la condition};
  \node [amil,below of =eval,fill=orange!20] (selection) {saute si elle est fausse};
  \uncover<4->{%
    \node [amil,below of =selection,fill=green!20] (bloc1) {bloc1};
    \node [amil,below of =bloc1,fill=orange!20] (saut1) {saut};
  }
  \node [amil,below of =saut1,fill=red!20](bloc2) {bloc2};
  \node [amil,below of =bloc2] (apres) {\emph{après}};
  \path [line] (selection) -| ++(2.6,0) |-  (bloc2);
\uncover<4->{\path [line] (saut1) -| ++(2.4,0) |-  (apres);}
\end{tikzpicture}
\end{column}
\end{columns}
\end{frame}

\mode<beamer>{\section[Démos]{Démos et fin}
\begin{frame}
  \frametitle{Démos et fin}
\end{frame}}


\section[Références]{Liens utiles}
\begin{frame}
\frametitle{Liens utiles\nowrite}

  \begin{itemize}
\item ma page : \url{http://www-lipn.univ-paris13.fr/~boudes/}
\item sites pour apprendre à développer :
  \begin{itemize}
  \item \url{http://www.siteduzero.com/} (chercher langage C)
  \item \url{http://www.developpez.com/} (chercher langage C)
  \end{itemize}
\item Le nouveau cours de terminal S :
  \url{https://science-info-lycee.fr}
\item Un livre de la BU: \emph{Le livre du C, premier langage (pour les
  vrais débutants en programmation)}, Claude Delannoy.
  \end{itemize}
\end{frame}

\begin{frame}
\frametitle{Liens utiles\nowrite}
\begin{itemize}
\item Polycopiés :
  \begin{itemize}
  \item le cours I3 de Laure Petrucci :\\
    {\small\url{http://www-lipn.univ-paris13.fr/~petrucci/polys.html}}
  \item le cours de Anne Canteaut :\\
    {\small\url{http://www-roc.inria.fr/secret/Anne.Canteaut/COURS_C/}}
  \item le cours de Bernard Cassagne :\\
    {\small\url{http://clips.imag.fr/commun/bernard.cassagne/Introduction_ANSI_C.html}}
  \item le cours de Henri Garreta : \\
    {\small\url{http://www.dil.univ-mrs.fr/~garreta/generique/}}
  \end{itemize}
\item Logiciels
  \begin{itemize}
  \item codeblocks : \url{http://www.codeblocks.org/}
  \item ubuntu : \url{http://www.ubuntu-fr.org/}
  \item virtualbox : \url{http://www.virtualbox.org/}
  \end{itemize}
  \end{itemize}
\end{frame}

\end{document}
